\section*{PHẦN MỞ ĐẦU}
\addcontentsline{toc}{section}{PHẦN MỞ ĐẦU}

\subsection*{1. Lý do chọn đề tài}
\addcontentsline{toc}{subsection}{1. Lý do chọn đề tài}

Trong hệ thống lý luận của chủ nghĩa duy vật lịch sử, mối quan hệ biện chứng giữa tồn tại xã hội và ý thức xã hội, giữa cơ sở hạ tầng và kiến trúc thượng tầng là nguyên lý nền tảng để nhận thức sự vận động của lịch sử. Triết học Mác -- Lênin chỉ rõ phương thức sản xuất vật chất quyết định đời sống tinh thần của xã hội; ngược lại, ý thức xã hội có tính độc lập tương đối và tác động trở lại thực tiễn thông qua hoạt động của con người. Trong xã hội trung đại Việt Nam, nho sĩ là tầng lớp trực tiếp tiếp nhận, truyền bá và thực thi hệ tư tưởng chính thống, giữ vai trò định hướng chuẩn mực chính trị, văn hóa và đạo đức quốc gia.

Thế kỷ XVIII đánh dấu bước ngoặt lớn của chế độ phong kiến Việt Nam khi thiết chế quân chủ lâm vào khủng hoảng ở cả Đàng Ngoài lẫn Đàng Trong, kéo theo sự suy giảm uy tín của Nho giáo chính thống vốn xác lập vị thế độc tôn từ thời Lê sơ. Cùng thời điểm đó, phong trào nông dân khởi nghĩa bùng nổ rộng khắp và đạt tới đỉnh cao với cuộc khởi nghĩa Tây Sơn do Quang Trung -- Nguyễn Huệ lãnh đạo, xóa bỏ thế phân liệt đất nước và đánh tan các thế lực ngoại xâm Xiêm, Thanh. Trước sự biến đổi căn bản của tồn tại xã hội, tầng lớp nho sĩ không còn giữ được sự thuần nhất về tư tưởng mà phân hóa mạnh mẽ về thái độ chính trị, đồng thời chuyển dịch từ lối học kinh viện Tống Nho sang khuynh hướng thực học, khảo chứng bách khoa và tinh thần nhân đạo.

Từ thực tế lịch sử đó, việc nghiên cứu đề tài ``Tìm hiểu về tầng lớp nho sĩ ở Việt Nam thế kỷ XVIII'' đặt ra yêu cầu thiết thực trên cả hai bình diện lý luận và thực tiễn. Về mặt lý luận, đề tài làm sáng tỏ cơ chế tác động của điều kiện kinh tế -- xã hội đối với sự biến đổi hệ tư tưởng, qua đó nhận diện vị trí trung gian và quy luật phân hóa của trí thức phong kiến trong các giai đoạn chuyển giao lịch sử. Về mặt thực tiễn, những kinh nghiệm lịch sử của kẻ sĩ thế kỷ XVIII gợi mở nhiều giá trị tham chiếu cho công tác xây dựng đội ngũ trí thức Việt Nam hiện nay: gắn kết lý luận với thực tiễn sản xuất, khắc phục bệnh chuộng hư danh và bằng cấp hình thức, bồi dưỡng bản lĩnh chính trị của người trí thức và hoàn thiện chính sách trọng dụng hiền tài phục vụ phát triển đất nước.

\subsection*{2. Tổng quan tình hình nghiên cứu}
\addcontentsline{toc}{subsection}{2. Tổng quan tình hình nghiên cứu}

Diện mạo lịch sử và tư tưởng của tầng lớp nho sĩ Việt Nam thế kỷ XVIII đã được giới nghiên cứu triết học, sử học và văn hóa học tiếp cận dưới ba nhóm công trình chính:

Thứ nhất là nhóm công trình nghiên cứu lịch sử tư tưởng và triết học Việt Nam. Tiêu biểu có bộ sách \textit{Lịch sử tư tưởng Việt Nam} (Tập 1, 1993) do Giáo sư, Tiến sĩ Nguyễn Tài Thư chủ biên đã khảo sát quá trình suy thoái của Nho giáo chính thống từ thế kỷ XVI đến thế kỷ XVIII và sự nảy nở của các yếu tố thực học. Trước đó, công trình \textit{Sự phát triển của tư tưởng ở Việt Nam từ thế kỷ XIX đến Cách mạng Tháng Tám} (Tập 1, 1973) của Giáo sư Trần Văn Giàu cũng phân tích cuộc khủng hoảng ý thức hệ phong kiến cuối thế kỷ XVIII như một tiền đề tư tưởng dẫn tới những biến chuyển ở giai đoạn cận đại.

Thứ hai là nhóm công trình lịch sử chế độ phong kiến và văn hóa giáo dục trung đại. Sách \textit{Lịch sử chế độ phong kiến Việt Nam} (Tập III, 1965) của các tác giả Phan Huy Lê, Chu Thiên, Vương Hoàng Tuyên, Đinh Xuân Lâm đã dựng lại bối cảnh mâu thuẫn xã hội, sự rệu rã của bộ máy quan trường và phong trào nông dân thế kỷ XVIII. Ở khía cạnh chế độ khoa cử, \textit{Việt Nam văn hóa sử cương} của Đào Duy Anh và \textit{Đại cương lịch sử văn hóa Việt Nam} (Tập 4, 1998) của Nguyễn Khắc Thuần cung cấp nhiều tư liệu xác thực về quy chế trường thi cùng những tiêu cực thời Lê mạt.

Thứ ba là nhóm chuyên khảo và bài báo khoa học về từng danh nhân, trào lưu học thuật cụ thể. Bài viết của Nguyễn Thị Hồng Phương trên \textit{Tạp chí Nghiên cứu Triết học} (2026) đã phân tích nhân sinh quan nhập thế gắn với quan niệm Lý -- Khí của Lê Quý Đôn; đồng thời bản dịch và khảo cứu tác phẩm \textit{Trúc Lâm tông chỉ nguyên thanh} (2001) làm rõ xu hướng dung hợp Tam giáo của Ngô Thì Nhậm. Cùng với các khảo cứu văn bản học của Nguyễn Kim Sơn (2004) về khuynh hướng khảo chứng học, các nghiên cứu của Đặng Việt Bích (2015) về thái độ của kẻ sĩ Bắc Hà đối với vua Quang Trung, của Quách Hiền (2018) về mạng lưới tư học làng xã, của Cao Văn Thức (2016) về học vị Sinh đồ và chuyên đề tư liệu về danh sư Võ Trường Toản (2020) đã bổ sung nguồn sử liệu phong phú trên cả hai miền Nam -- Bắc.

Kế thừa kết quả của các công trình đi trước, tiểu luận hệ thống hóa bức tranh toàn cảnh về tầng lớp nho sĩ thế kỷ XVIII dưới phương pháp luận duy vật biện chứng và duy vật lịch sử, gắn kết sự phân hóa chính trị với bước chuyển học thuật để rút ra bài học đối với đội ngũ trí thức đương đại.

\subsection*{3. Mục tiêu và nhiệm vụ nghiên cứu}
\addcontentsline{toc}{subsection}{3. Mục tiêu và nhiệm vụ nghiên cứu}

Mục tiêu tổng quát của đề tài là vận dụng lý luận triết học Mác -- Lênin để làm rõ nguồn gốc xã hội, thực trạng phân hóa chính trị và sự vận động tư tưởng -- học thuật của tầng lớp nho sĩ Việt Nam thế kỷ XVIII; trên cơ sở đó đánh giá khách quan những đóng góp, hạn chế lịch sử và rút ra bài học đối với việc phát huy vai trò đội ngũ trí thức Việt Nam hiện nay.

Để thực hiện mục tiêu trên, đề tài giải quyết ba nhiệm vụ cụ thể tương ứng với ba chương:

\textit{Nhiệm vụ thứ nhất,} xác lập cơ sở lý luận duy vật lịch sử về ý thức xã hội và đặc trưng của tầng lớp trí thức; phân tích vị thế nho sĩ trước thế kỷ XVIII cùng những biến động kinh tế, chính trị, khoa cử thế kỷ XVIII tác động đến tầng lớp này.

\textit{Nhiệm vụ thứ hai,} phân tích các khuynh hướng ứng xử chính trị của nho sĩ trước biến cố thời cuộc và làm sáng tỏ những chuyển biến về thế giới quan triết học, phương pháp khảo chứng, tinh thần dung hợp tôn giáo và giá trị nhân văn trong sáng tác của họ.

\textit{Nhiệm vụ thứ ba,} đánh giá giá trị tiến bộ và giới hạn lịch sử của nho sĩ thế kỷ XVIII theo nguyên tắc lịch sử -- cụ thể, từ đó liên hệ rút ra các bài học thực tiễn trong công tác xây dựng đội ngũ trí thức của Đảng và Nhà nước hiện nay.

\subsection*{4. Đối tượng và phạm vi nghiên cứu}
\addcontentsline{toc}{subsection}{4. Đối tượng và phạm vi nghiên cứu}

Đối tượng nghiên cứu của đề tài là tầng lớp nho sĩ Việt Nam thế kỷ XVIII, tập trung vào sự phân hóa lập trường chính trị -- xã hội và sự chuyển biến trong thế giới quan triết học, nhân sinh quan, khuynh hướng học thuật dưới sự quy định của hoàn cảnh lịch sử đương thời.

Về phạm vi không gian, đề tài khảo sát hoạt động và tư tưởng của nho sĩ trên cả hai vùng lãnh thổ Đàng Ngoài và Đàng Trong của nước Đại Việt. Về phạm vi thời gian, nội dung nghiên cứu giới hạn trọng tâm trong thế kỷ XVIII (từ thời điểm khủng hoảng của chính quyền Lê -- Trịnh, chúa Nguyễn đến giai đoạn triều đại Tây Sơn), có đối chiếu với tiền đề Nho giáo các thế kỷ trước và liên hệ với thực tiễn Việt Nam hiện nay.

\subsection*{5. Phương pháp nghiên cứu}
\addcontentsline{toc}{subsection}{5. Phương pháp nghiên cứu}

Đề tài này được thực hiện trên cơ sở phương pháp luận của chủ nghĩa duy vật biện chứng và chủ nghĩa duy vật lịch sử. Ngoài ra, đề tài còn sử dụng các phương pháp khác như:

Phương pháp thống nhất lịch sử và logic giúp tái hiện diễn biến lịch sử thế kỷ XVIII theo trình tự thời gian, qua đó nắm bắt quy luật vận động nội tại của tư tưởng nho sĩ mà không sa vào lối liệt kê sự kiện thuần túy.

Phương pháp phân tích và tổng hợp được áp dụng để bóc tách từng thái độ xuất xử, từng luận điểm triết học của các học giả tiêu biểu, từ đó khái quát thành diện mạo chung của trí thức thời kỳ này.

Phương pháp so sánh, đối chiếu giúp nhận diện sự khác biệt giữa bộ phận nho sĩ bảo thủ và bộ phận nho sĩ cấp tiến, giữa đặc trưng kẻ sĩ Bắc Hà với sĩ phu Gia Định, cũng như giữa Lý học Tống Nho kinh viện với trào lưu thực học.

Phương pháp tiếp cận hệ thống đặt tầng lớp nho sĩ trong mối liên hệ chỉnh thể với phương thức sản xuất phong kiến, thiết chế khoa cử và phong trào nông dân khởi nghĩa thế kỷ XVIII.

\subsection*{6. Kết cấu của đề tài}
\addcontentsline{toc}{subsection}{6. Kết cấu của đề tài}

Ngoài phần mở đầu, kết luận, tài liệu tham khảo, nội dung của đề tài này bao gồm 03 Chương, 06 Tiết. Mạch nghiên cứu đi từ nền tảng lý luận và điều kiện lịch sử khách quan ở Chương 1, sang phân tích thực trạng phân hóa chính trị cùng bước chuyển tư tưởng ở Chương 2, trước khi tổng kết giá trị, hạn chế và đúc rút bài học thực tiễn tại Chương 3.
